\section{Discussion}
\label{sec:discussion}
The runtime makes the intended agent-mediated sequence explicit: the agent acquires configured process observations, analyzes them against an objective, submits a control intent, receives a command outcome or supported readback, and can request later DAQ observations for a subsequent decision. The selected nominal fixed-horizon batch had lower retained-window accumulated error with the deterministic policy than with no-write in all nine selected valid blocks. An earlier valid block was repeated in the selected batch, and the available records do not explain the batch-selection rule. There were only three selected blocks per scene. Initial process values were constrained to windows rather than required to match exactly, and no named conventional controller was tested. The result therefore supports a comparison specific to these simulations, not a claim of general process-control efficacy.

Fig.~\ref{fig:optimization-trajectory} presents the recorded injection mass traces and pressure and hold-time requests alongside the selected retained-window IAE contrasts. These records do not establish the causal effect of each setting: iteration-level timing and independent per-request server records are unavailable.

The mapping and guard tests support only the configured branches and inputs that were exercised. The test oracles share configuration with the tested limits, and the mock-driver ablation lacks independent dispatch records. Simulator readback is not physical actuation evidence. The LLM-driven episodes exercised the agent-facing recipe tool and returned acknowledged updates, but all ended below target. Without the server journal, those responses cannot establish a complete agent-to-device write chain. Agent task completion, command acknowledgement, and process attainment must therefore remain separate outcomes.

The platform and simulator were already running, and a complete immutable service and configuration snapshot was not retained. Application checks complement PLC, safety-instrumented-system, and equipment protections. Further process validation requires a larger pre-specified study with verified reset states, common sampling and horizons, a named controller baseline, and complete trajectories. Agent evaluation requires a fixed acquisition-and-control task, retained model and runtime versions, and correlated observation-to-intent-to-device records for every episode. Human-approval effects and physical safety need separate studies.

\section{Conclusion}
\label{sec:conclusion}
This work presents a configuration-driven runtime for agent-mediated online data acquisition and governed setpoint control. The agent-facing path binds production context and direction-specific acquisition/control signals, returns observations for objective-conditioned analysis, and carries a typed control intent through route-dependent checks before dispatch while keeping command outcomes distinct from readback and later process observations. Selected mapping and guard branches were exercised across simulated scenes and mock-driver fixtures. In the selected nominal fixed-horizon simulator batch, all nine valid scene--seed pairs had lower retained-window IAE under the deterministic setpoint policy than under no-write; an earlier valid pair was repeated and the batch-selection rule is unresolved. This exploratory result lacks a conventional controller baseline. The injection traces show recorded setpoint requests and process observations; three LLM-driven episodes exercised the agent-facing recipe tool but were canceled below target and lacked a complete server-side write record. The evidence establishes an application-level agent-to-operation integration boundary, not general optimizer convergence, autonomous-agent efficacy, a complete agent-to-device trace, or physical-plant safety.
